\documentclass[10pt,a4paper]{extarticle}
\usepackage[utf8]{inputenc}
\usepackage[T1]{fontenc}
\usepackage[french]{babel}
\usepackage{geometry}
\geometry{hmargin=0.85cm,vmargin=0.75cm}
\usepackage{graphicx}
\usepackage{float}
\usepackage{microtype}

\setlength{\parindent}{0pt}
\setlength{\parskip}{2pt}
\renewcommand{\baselinestretch}{0.96}

% Captions plus grandes et plus lisibles
\usepackage[font=small,labelfont=bf,skip=3pt]{caption}

\begin{document}

\section*{Data visualisation : lecture des facteurs associés au diabète}


Dans cette partie, nous ne cherchons pas à afficher toutes les variables du jeu de données, mais à sélectionner les visualisations qui font apparaître les tendances les plus lisibles. L'objectif est de passer d'une lecture brute des colonnes à une lecture par profils, en montrant comment le BMI, l'âge, l'état de santé et certains indicateurs médicaux structurent la présence du diabète dans les données.

\begin{figure}[H]
    \centering
    \includegraphics[width=0.42\linewidth]{Figure_3.png}
    \caption{Distribution du BMI selon le statut diabétique. Les patients diabétiques présentent globalement des valeurs de BMI plus élevées, mais les deux distributions se chevauchent nettement.}
    \label{fig:bmi_distribution}
\end{figure}

La figure \ref{fig:bmi_distribution} nuance la lecture du poids : elle montre la dispersion complète plutôt qu'une simple moyenne. Le décalage du BMI vers le haut chez les patients diabétiques est visible, mais le chevauchement important rappelle que le diabète ne se résume pas au poids. Le BMI est donc un signal important, lié au mode de vie, mais il doit être interprété avec d'autres dimensions du profil patient.

\begin{figure}[H]
    \centering
    \includegraphics[width=0.94\linewidth]{Figure_4.png}
    \caption{Carte de risque : taux de diabète selon le groupe d'âge et la catégorie de BMI.}
    \label{fig:age_bmi_heatmap}
\end{figure}

La figure \ref{fig:age_bmi_heatmap} est la visualisation la plus structurante. Elle croise une variable non modifiable, l'âge, avec un facteur fortement lié au style de vie, le BMI. Chez les individus avec un BMI normal, le taux de diabète reste inférieur à 12\% même aux âges élevés. Chez les patients obèses, le taux dépasse 20\% dès 50--54 ans, atteint 29,2\% entre 60 et 64 ans, puis culmine autour de 35\% entre 70 et 74 ans. Cette carte montre que le risque dépend d'un profil : l'âge amplifie fortement l'effet du BMI. Le message métier est clair : on ne peut pas agir sur l'âge, mais le contrôle du poids peut limiter l'entrée dans les zones de risque les plus élevées.

\begin{figure}[H]
    \centering
    \begin{minipage}[t]{0.485\linewidth}
        \centering
        \includegraphics[width=\linewidth]{Figure_5.png}
    \end{minipage}
    \hfill
    \begin{minipage}[t]{0.485\linewidth}
        \centering
        \includegraphics[width=\linewidth]{Figure_6.png}
    \end{minipage}
    \caption{Cartes de chaleur des signaux cliniques et comportementaux : hypertension/cholestérol à gauche, santé générale/activité physique à droite.}
    \label{fig:clinical_behavioral_heatmaps}
\end{figure}

La carte hypertension--cholestérol montre une hausse progressive du taux de diabète lorsque les signaux médicaux se cumulent : 4,19\% sans facteur, 10,42\% avec cholestérol seul, 16,73\% avec hypertension seule, et 29,71\% lorsque les deux sont présents. La carte santé générale--activité physique relie un indicateur subjectif de santé à un comportement modifiable : à chaque niveau de santé déclarée, les personnes actives présentent un taux plus faible que les personnes inactives. L'activité physique apparaît donc comme un facteur associé à un profil plus favorable, sans annuler l'effet d'une mauvaise santé générale.





\begin{figure}[H]
    \centering
    \includegraphics[width=0.54\linewidth]{Figure_8.png}
    \caption{Corrélation des variables avec \texttt{Diabetes\_binary}. Cette lecture synthétique met en évidence les variables les plus associées à la cible.}
    \label{fig:corr_target}
\end{figure}

La figure \ref{fig:corr_target} rend la matrice précédente plus lisible en isolant la corrélation avec la cible. Elle confirme le poids des variables de santé générale, d'hypertension, de difficulté à marcher, de BMI, de cholestérol et d'âge. Les corrélations négatives de \texttt{PhysActivity}, \texttt{Education} et \texttt{Income} suggèrent des profils statistiquement plus favorables, mais elles doivent rester interprétées comme des associations. Le résultat important n'est donc pas l'existence d'un facteur unique, mais la convergence de plusieurs signaux cohérents.

\vspace{0.35cm}

\textbf{Synthèse.} L'analyse visuelle fait ressortir une logique d'accumulation du risque. Le BMI est central, mais il devient beaucoup plus informatif lorsqu'il est croisé avec l'âge. L'hypertension et le cholestérol servent de signaux d'alerte métaboliques, tandis que l'activité physique est associée à des taux plus faibles, y compris à santé générale comparable. Ces résultats ne démontrent pas une causalité directe, mais ils orientent clairement la suite de l'étude vers le feature engineering et la régression logistique.

\end{document}


